% ============================================================
%  CHAPTER 12: SURFACE TENSION, VISCOSITY AND CAPILLARITY
%  The Physics Codex: A Complete University Foundation
%  Korede Samuel Omotosho (Falcon)
%  Aurikrex Learning Systems, 2026
%
%  File: chapters/part2/ch12_surface_tension.tex
% ============================================================

\chapter{Surface Tension, Viscosity and Capillarity}
\label{ch:surface_tension}
\minitoc
\newpage

\chapterquote{In every walk with nature, one receives
far more than he seeks.}{John Muir}

\begin{objectivebox}
By the end of this chapter, you should be able to:
\begin{enumerate}[noitemsep]
  \item Explain surface tension in terms of
  intermolecular forces
  \item Define surface tension and calculate it
  using the force per unit length definition
  \item Explain and calculate the excess pressure
  inside a soap bubble and a liquid drop
  \item Define capillarity and derive the formula
  for capillary rise
  \item Explain why water rises in glass capillaries
  but mercury falls
  \item Define viscosity and explain its molecular
  origin
  \item State Stokes' law and use it to calculate
  terminal velocity in a viscous fluid
  \item Distinguish between laminar and turbulent
  flow
\end{enumerate}
\end{objectivebox}

% ============================================================
\section{Intuition: The Invisible Skin}
% ============================================================

Fill a glass of water very carefully until it is
slightly above the rim. Instead of immediately
spilling over, the water forms a gentle dome above
the glass edge --- held there by something. A small
insect called a water strider walks across the surface
of a pond as if the water were solid ground, its legs
making small dimples but never breaking through.
A steel needle, denser than water, can be made to
float if placed gently and horizontally on the surface.

None of these phenomena makes sense if we think of
a liquid surface as simply nothing --- as just
the boundary between water and air. Something is
acting there. That something is \textbf{surface tension}
--- an emergent property of the intermolecular forces
that hold a liquid together, creating what behaves
like an elastic membrane stretched over the surface.

% ============================================================
\section{Intermolecular Forces: The Origin of Surface
Tension}
% ============================================================

To understand surface tension, we must first
understand where it comes from.

Inside a liquid, each molecule is surrounded on
all sides by neighbouring molecules. The attractive
forces (\textbf{cohesive forces}) pull equally in
all directions. The net force on an interior molecule
is zero.

At the surface, the situation is different. A surface
molecule has neighbours below and to its sides, but
very few neighbours above (only air molecules, which
are much farther apart). The cohesive forces pulling
downward and sideways are not balanced by forces from
above. The surface molecule experiences a net
\textbf{inward} force.

\begin{diagbox}[Interior vs Surface Molecule: Intermolecular Forces]
\begin{center}
\begin{tikzpicture}[scale=1.0, font=\small]

  % Liquid region
  \fill[codexblue!15] (0,0) rectangle (8,4);
  % Air region
  \fill[cyan!5] (0,4) rectangle (8,6);
  % Surface line
  \draw[codexblue, very thick] (0,4) -- (8,4)
    node[right]{Surface};
  % Air label
  \node[gray, font=\itshape] at (3.3,5.2) {Air (few molecules, weak forces)};
  % Liquid label
  \node[codexblue, font=\itshape] at (5.5,0.5) {Liquid interior};

  % ---- INTERIOR MOLECULE ----
  \fill[codexred] (2.5,2.0) circle (8pt)
    node[below=13pt, font=\bfseries\small, black]{Interior molecule};

  % Force arrows in all directions (equal)
  \foreach \angle in {0,45,90,135,180,225,270,315} {
    \draw[->, thick, codexred!70]
      ({2.5 + 0.15*cos(\angle)}, {2.0 + 0.15*sin(\angle)})
      -- ({2.5 + 0.7*cos(\angle)}, {2.0 + 0.7*sin(\angle)});
  }
  % Net force = zero
  \node[codexgreen, font=\small] at (2.5,2.0)
    {};
  \draw[codexgreen, thick] (2.5,0.9)
    node[right, font=\tiny]{Net force $= 0$};

  % ---- SURFACE MOLECULE ----
  \fill[codexgold] (5.5,4.0) circle (8pt)
    node[above=10pt, font=\bfseries\small, black]{Surface molecule};

  % Force arrows only downward and sideways (no upward)
  \foreach \angle in {0,180,225,270,315} {
    \draw[->, thick, codexgold!80]
      ({5.5 + 0.15*cos(\angle)}, {4.0 + 0.15*sin(\angle)})
      -- ({5.5 + 0.7*cos(\angle)}, {4.0 + 0.7*sin(\angle)});
  }
  % Very weak upward forces from air
  \foreach \angle in {45,90,135} {
    \draw[->, thin, gray, dashed]
      ({5.5 + 0.15*cos(\angle)}, {4.0 + 0.15*sin(\angle)})
      -- ({5.5 + 0.35*cos(\angle)}, {4.0 + 0.35*sin(\angle)});
  }
  \node[gray, font=\tiny, align=center] at (7.1,5.2) {Weak air\\forces};

  % Net inward force arrow
  \draw[->, very thick, codexred, line width=2pt]
    (5.5,3.9) -- (5.5,3.0)
    node[right]{Net inward force};

  % Consequence annotation
  \node[align=center, font=\small] at (4.0,-0.9)
    {Surface molecules are pulled inward\\
    $\implies$ Surface tends to contract\\
    $\implies$ Acts like a stretched elastic membrane};

\end{tikzpicture}
\end{center}
\end{diagbox}

This net inward force has two consequences. First,
molecules at the surface have higher potential energy
than those inside (work must be done to bring a
molecule from inside to the surface). Second, the
surface tends to \textit{contract} --- to minimise
the number of molecules at the surface, and therefore
minimise the total surface area. This tendency to
contract is what we experience as \textbf{surface
tension}.

% ============================================================
\section{Defining Surface Tension}
% ============================================================

\begin{definitionbox}[Surface Tension]
\textbf{Surface tension} $\gamma$ (also written $T$)
is the force per unit length acting along the surface
of a liquid, perpendicular to any imaginary line
drawn on the surface:
\[
  \gamma = \frac{F}{L}
\]
where $F$ is the force acting along a length $L$ of
the surface. SI unit: N\,m$^{-1}$.

Equivalently, surface tension is the energy per unit
area of the surface:
\[
  \gamma = \frac{W}{A}\ \text{(J\,m}^{-2}\text{)}
\]
Both definitions give the same numerical value and
are physically equivalent.
\end{definitionbox}

\begin{table}[H]
\centering
\caption{Surface Tension of Common Liquids at 20°C}
\label{tab:surface_tension}
\begin{tabular}{ll}
\toprule
\textbf{Liquid} & \textbf{Surface Tension (N\,m$^{-1}$)}\\
\midrule
Mercury & 0.480\\
Water & 0.073\\
Glycerol & 0.063\\
Olive oil & 0.033\\
Soapy water & 0.025--0.040\\
Ethanol & 0.022\\
Acetone & 0.024\\
\bottomrule
\end{tabular}
\end{table}

\begin{keybox}[Factors Affecting Surface Tension]
\begin{itemize}[noitemsep]
  \item \textbf{Temperature:} Surface tension
  \textit{decreases} as temperature increases
  (molecules have more energy, cohesive forces
  effectively weaker)
  \item \textbf{Contamination:} Soap and detergents
  dramatically \textit{reduce} surface tension (they
  are surfactants --- surface-active agents)
  \item \textbf{Type of liquid:} Depends on the
  nature and strength of intermolecular forces
\end{itemize}
\end{keybox}

\begin{examplebox}[12.1]
\stepcounter{workedexample}
A thin wire frame with a movable crossbar is used to
measure surface tension. The crossbar has length
$4.0\ \text{cm}$. A force of $5.84\times10^{-3}\ \text{N}$
is required to slowly pull the crossbar away from
the liquid film.\\
(a) Explain why the factor of 2 appears in the
calculation.\\
(b) Calculate the surface tension.
\end{examplebox}

\begin{solutionbox}
\textbf{(a)} A liquid film has \textit{two} surfaces
--- a front face and a back face. The force must
stretch both surfaces simultaneously, so the
effective length resisting the force is $2L$.

\textbf{(b)}
\[
  \gamma = \frac{F}{2L}
  = \frac{5.84\times10^{-3}}{2 \times 0.040}
  = \frac{5.84\times10^{-3}}{0.080}
  = 0.073\ \text{N\,m}^{-1}
\]
This is consistent with the surface tension of
water at room temperature.
\end{solutionbox}

% ============================================================
\section{Excess Pressure in Bubbles and Drops}
% ============================================================

Surface tension creates a pressure difference across
any curved liquid surface --- the pressure inside is
always greater than the pressure outside.

\begin{processbox}[Excess Pressure Inside a Liquid Drop]
Consider a spherical liquid drop of radius $r$.
Surface tension acts along the equator trying to
contract the sphere. If we imagine cutting the drop
in half, the surface tension force pulls the two
halves together along the circumference $2\pi r$:
\[
  F_{surface} = \gamma \times 2\pi r
\]
The excess pressure $\Delta P$ inside acts over the
circular cross-section:
\[
  F_{pressure} = \Delta P \times \pi r^2
\]
At equilibrium: $\Delta P \times \pi r^2
= \gamma \times 2\pi r$
\[
  \boxed{\Delta P = \frac{2\gamma}{r}}
  \quad \text{(liquid drop --- one surface)}
\]
\end{processbox}

\begin{processbox}[Excess Pressure Inside a Soap Bubble]
A soap bubble has \textbf{two} surfaces (inner and
outer). Both contribute to the contracting force:
\[
  F_{surface} = 2 \times \gamma \times 2\pi r
  = 4\pi r \gamma
\]
\[
  \Delta P \times \pi r^2 = 4\pi r \gamma
\]
\[
  \boxed{\Delta P = \frac{4\gamma}{r}}
  \quad \text{(soap bubble --- two surfaces)}
\]
\end{processbox}

\begin{diagbox}[Excess Pressure: Drop vs Bubble]
\begin{center}
\begin{tikzpicture}[scale=1.0, font=\small]

  % === LIQUID DROP ===
  \begin{scope}[shift={(0,0)}]
    \node[font=\bfseries] at (2,4.8) {Liquid Drop};
    % Drop
    \fill[codexblue!40] (2,2) circle (1.8cm);
    \draw[very thick, codexblue] (2,2) circle (1.8cm);
    % One surface label
    \node[codexblue, font=\small] at (2.0,4.3)
      {One surface};
    \draw[->, codexblue] (3.3,3.3) -- (2.8,2.8);
    % Radius
    \draw[<->, gray] (2,2) -- (3.27,0.73)
      node[midway, below right]{$r$};
    % Excess pressure is normal to the interface; surface tension is tangential.
    \draw[->, thick, codexred] (3.8,2.0) -- (4.35,2.0);
    \node[codexred, font=\tiny, align=center] at (4.6,2.9)
      {$\Delta P$: pressure\\normal to surface};
    \draw[<->, thick, codexgreen] (3.8,1.45) -- (3.8,2.55);
    \node[codexgreen, font=\tiny, align=center] at (4.6,1.25)
      {Surface tension $\gamma$: \\tangent to surface};
    % Excess pressure label
\node[font=\small\bfseries, codexblue] at (2,-0.4)
      {$\Delta P = \dfrac{2\gamma}{r}$};
    % Higher pressure inside
    \node[font=\tiny, align=center] at (2,-1.5)
      {$P_{inside} > P_{outside}$\\
      by $\dfrac{2\gamma}{r}$};
  \end{scope}

  % === SOAP BUBBLE ===
  \begin{scope}[shift={(6,0)}]
    \node[font=\bfseries] at (2,4.8) {Soap Bubble};
    % Film between surfaces
    \fill[codexblue!20] (2,2) circle (2.0cm);
    \fill[white] (2,2) circle (1.6cm);
    % Surface outlines are drawn after the fills to remain visible
    \draw[very thick, codexblue] (2,2) circle (2.0cm);
    \draw[very thick, codexblue!60] (2,2) circle (1.6cm);
    % Labels
    \node[codexblue, font=\tiny] at (3.5,3.6)
      {Outer surface};
    \draw[->, codexblue, thin] (3.3,3.4)
      -- (2.9,2.9);
    \node[codexblue!70, font=\tiny] at (3.4,2.8)
      {Inner surface};
    \draw[->, codexblue!70, thin] (3.2,2.7)
      -- (2.7,2.4);
    % Two surfaces annotation
    \node[codexblue, font=\small] at (2,4.3)
      {Two surfaces};
    % Air inside label
    \node[font=\tiny] at (2,2) {Air};
    % Radius
    \draw[<->, codexred] (2,2) -- (4.0,2)
      node[midway, above]{$r$};
    % Excess pressure
    \node[font=\small\bfseries, codexblue] at (2,-0.5)
      {$\Delta P = \dfrac{4\gamma}{r}$};
    \node[font=\tiny, align=center] at (2,-1.5)
      {Twice the excess pressure\\
      of a liquid drop};
  \end{scope}

\end{tikzpicture}
\end{center}
\end{diagbox}

\begin{examplebox}[12.2]
\stepcounter{workedexample}
A soap bubble has radius $3.0\ \text{cm}$.
The surface tension of the soap solution is
$0.030\ \text{N\,m}^{-1}$.\\
(a) Calculate the excess pressure inside the bubble.\\
(b) What is the total pressure inside the bubble if
atmospheric pressure is $1.013\times10^5\ \text{Pa}$?\\
(c) Two soap bubbles of radii $r_1 = 2.0\ \text{cm}$
and $r_2 = 4.0\ \text{cm}$ are joined together.
Which direction does air flow when the junction is
opened? Explain.
\end{examplebox}

\begin{solutionbox}
\textbf{(a)} Excess pressure (soap bubble, two surfaces):
\[
  \Delta P = \frac{4\gamma}{r}
  = \frac{4 \times 0.030}{0.030}
  = \frac{0.120}{0.030} = 4.0\ \text{Pa}
\]

\textbf{(b)} Total pressure:
\[
  P_{total} = P_{atm} + \Delta P
  = 1.013\times10^5 + 4.0
  \approx 1.013\times10^5\ \text{Pa}
\]
The excess pressure is tiny compared to atmospheric
pressure --- which is why blowing a soap bubble
requires very little extra effort.

\textbf{(c)} Excess pressure in small bubble:
$\Delta P_1 = 4\gamma/r_1 = 4(0.030)/0.020
= 6.0\ \text{Pa}$

Excess pressure in large bubble:
$\Delta P_2 = 4\gamma/r_2 = 4(0.030)/0.040
= 3.0\ \text{Pa}$

The small bubble has higher internal pressure.
When joined, air flows \textbf{from the small bubble
into the large bubble}. The small bubble shrinks and
the large one grows. This is counterintuitive ---
one might expect the large bubble to push into the
small one. The physics says otherwise: smaller radius
$\implies$ higher pressure $\implies$ air flows outward.
\end{solutionbox}

% ============================================================
\section{Capillarity}
% ============================================================

\begin{definitionbox}[Capillarity]
\textbf{Capillarity} (or capillary action) is the
spontaneous rise or fall of a liquid in a narrow tube
(capillary tube) due to the combined effects of
surface tension (cohesion within the liquid) and
adhesion between the liquid and the tube wall.
\end{definitionbox}

Whether a liquid rises or falls in a capillary tube
depends on whether adhesion or cohesion dominates:

\begin{itemize}
  \item \textbf{Water in glass:} Adhesion of water
  to glass $>$ cohesion of water to water. Water
  \textit{wets} glass. The meniscus is \textit{concave}.
  Water \textbf{rises} in the tube.

  \item \textbf{Mercury in glass:} Cohesion of
  mercury to mercury $>$ adhesion of mercury to glass.
  Mercury does \textit{not} wet glass. The meniscus
  is \textit{convex}. Mercury \textbf{falls} in the tube.
\end{itemize}

\begin{diagbox}[Capillary Rise and Fall: Water vs Mercury]
\begin{center}
\begin{tikzpicture}[scale=1.0, font=\small]
  % === WATER IN GLASS (rises) ===
  \begin{scope}[shift={(0,0)}]
    \node[font=\bfseries, align=center] at (1.5,7.0)
      {Water in Glass\\(rises)};
    % Outer container
    \fill[codexblue!10] (0,0) rectangle (3.0,4.5);
    \draw[thick] (0,0) rectangle (3.0,4.5);
    % Bulk water level
    \draw[codexblue, thick] (0,2.0) -- (0.8,2.0);
    \draw[codexblue, thick] (2.2,2.0) -- (3.0,2.0);
    \fill[codexblue!25] (0,0) rectangle (3.0,2.0);
    % Capillary tube
    \draw[thick] (0.9,0) -- (0.9,5.5);
    \draw[thick] (2.1,0) -- (2.1,5.5);
    % Water inside capillary (risen)
    \fill[codexblue!40] (0.9,0) rectangle (2.1,3.8);
    % Concave meniscus at top
    \draw[codexblue, very thick]
      (0.9,3.8) .. controls (1.1,3.4) and (1.9,3.4) ..
      (2.1,3.8);
    \fill[codexblue!40]
      (0.9,3.8) .. controls (1.1,3.4) and (1.9,3.4) ..
      (2.1,3.8) -- (2.1,3.8) -- (0.9,3.8);
    % Height h
    \draw[<->, codexred, thick]
      (2.3,2.0) -- (2.3,3.6);
    \node[right, codexred] at (2.4,2.8) {$h$};
    % Contact angle at the right wall: measured through the liquid.
    \draw[densely dashed, gray] (2.1,3.8) -- (1.75,3.1);
    \draw[codexred, thin]
      (2.1,3.45) arc[start angle=270,end angle=243,radius=0.35];
    \node[codexred, font=\tiny] at (1.88,3.35) {$\theta$};
    \node[font=\tiny, codexblue] at (1.5,5.95)
      {Concave meniscus};
    \node[font=\tiny, codexblue] at (1.5,6.35)
      {(adhesion $>$ cohesion)};
    % Arrows: adhesion pulls up (label moved to arrow midpoint,
    % clear of the meniscus labels above it)
    \draw[->, thick, codexgreen]
      (0.9,3.5) -- (0.9,4.2)
      node[midway, left, font=\tiny]{adhesion};
  \end{scope}
  % === MERCURY IN GLASS (falls) ===
  \begin{scope}[shift={(5.0,0)}]
    \node[font=\bfseries, align=center] at (1.5,7.0)
      {Mercury in Glass\\(falls)};
    % Outer container
    \fill[gray!15] (0,0) rectangle (3.0,4.5);
    \draw[thick] (0,0) rectangle (3.0,4.5);
    % Bulk mercury level
    \draw[gray!70, thick] (0,2.5) -- (0.8,2.5);
    \draw[gray!70, thick] (2.2,2.5) -- (3.0,2.5);
    \fill[gray!50] (0,0) rectangle (3.0,2.5);
    % Capillary tube
    \draw[thick] (0.9,0) -- (0.9,5.5);
    \draw[thick] (2.1,0) -- (2.1,5.5);
    % Mercury inside capillary (depressed)
    \fill[gray!60] (0.9,0) rectangle (2.1,1.8);
    % Convex meniscus at top (bulges upward)
    \draw[gray!80, very thick]
      (0.9,1.8) .. controls (1.1,2.3) and (1.9,2.3) ..
      (2.1,1.8);
    \fill[gray!60]
      (0.9,1.8) .. controls (1.1,2.3) and (1.9,2.3) ..
      (2.1,1.8) -- (2.1,1.8) -- (0.9,1.8);
    % Depression h
    \draw[<->, codexred, thick]
      (2.3,1.8) -- (2.3,2.5);
    \node[right, codexred] at (2.4,2.15)
      {\shortstack{$h$\\(depressed)}};
    % Contact angle at the right wall: measured through the liquid.
    \draw[densely dashed, gray] (2.1,1.8) -- (1.85,2.43);
    \draw[codexred, thin]
      (1.969,2.125) arc[start angle=112,end angle=270,radius=0.35];
    \node[codexred, font=\tiny] at (1.78,2.12) {$\theta$};
    % Move Mercury labels above the tube for contrast and to avoid glass-wall overlap.
    \node[font=\tiny, black] at (1.5,5.95)
      {Convex meniscus};
    \node[font=\tiny, black] at (1.5,6.35)
      {(cohesion $>$ adhesion)};
    % Arrows: cohesion pulls down
    \draw[->, thick, codexred]
      (1.5,1.8) -- (1.5,1.1)
      node[right, font=\tiny, align=left]{cohesion\\pulls down};
  \end{scope}
  \node[font=\tiny, align=center] at (4.0,-0.55)
    {Contact angle $\theta$ is measured through the liquid};
\end{tikzpicture}
\end{center}
\end{diagbox}

\begin{processbox}[Deriving the Capillary Rise Formula]
For a liquid of surface tension $\gamma$, density
$\rho$, in a tube of radius $r$, with contact angle
$\theta$:

The surface tension force acting upward around the
contact line:
\[
  F_{up} = \gamma \cos\theta \times 2\pi r
\]

Weight of liquid column raised to height $h$:
\[
  W = \rho g h \pi r^2
\]

At equilibrium ($F_{up} = W$):
\[
  \gamma \cos\theta \times 2\pi r = \rho g h \pi r^2
\]
\[
  \boxed{h = \frac{2\gamma\cos\theta}{\rho g r}}
\]
For water on glass, $\theta \approx 0°$,
so $\cos\theta \approx 1$:
\[
  h \approx \frac{2\gamma}{\rho g r}
\]
\textbf{Narrower tube $\implies$ greater rise.}
\end{processbox}

\begin{diagbox}[Capillary Rise Formula: Narrower Tube Rises Higher]
\begin{center}
\begin{tikzpicture}[scale=0.9, font=\small]

  % Container
  \fill[codexblue!10] (0,0) rectangle (9,2);
  \draw[thick] (0,0) rectangle (9,2);
  % Bulk water level
  \draw[codexblue, thick] (0,2) -- (9,2);
  \fill[codexblue!20] (0,0) rectangle (9,2);

  % Tube 1 (widest, r = large, lowest rise)
  \begin{scope}[shift={(1,0)}]
    \draw[thick] (0,0) -- (0,6); % left wall
    \draw[thick] (1.6,0) -- (1.6,6); % right wall
    \fill[codexblue!40] (0,0) rectangle (1.6,2.8);
    \draw[codexblue, very thick]
      (0,2.8) .. controls (0.25,2.4) and (1.35,2.4) ..
      (1.6,2.8);
    \draw[<->, codexred] (1.8,2.0) -- (1.8,2.4);
    \node[right, codexred, font=\tiny] at (1.8,2.2)
      {\shortstack{$h_1$}};
    \node[font=\tiny, align=center] at (0.8,5.5)
      {Wide\\$r_1$};
  \end{scope}

  % Tube 2 (medium)
  \begin{scope}[shift={(4,0)}]
    \draw[thick] (0,0) -- (0,6);
    \draw[thick] (1.0,0) -- (1.0,6);
    \fill[codexblue!40] (0,0) rectangle (1.0,3.5);
    \draw[codexblue, very thick]
      (0,3.5) .. controls (0.15,3.1) and (0.85,3.1) ..
      (1.0,3.5);
    \draw[<->, codexred] (1.2,2.0) -- (1.2,3.1);
    \node[right, codexred, font=\tiny] at (1.3,2.55)
      {$h_2$};
    \node[font=\tiny, align=center] at (0.5,5.5)
      {Medium\\$r_2$};
  \end{scope}

  % Tube 3 (narrowest, highest rise)
  \begin{scope}[shift={(6.5,0)}]
    \draw[thick] (0,0) -- (0,6);
    \draw[thick] (0.5,0) -- (0.5,6);
    \fill[codexblue!40] (0,0) rectangle (0.5,4.8);
    \draw[codexblue, very thick]
      (0,4.8) .. controls (0.06,4.4) and (0.44,4.4) ..
      (0.5,4.8);
    \draw[<->, codexred] (0.7,2.0) -- (0.7,4.4);
    \node[right, codexred, font=\tiny] at (0.78,3.2)
      {\shortstack{$h_3$}};
    \node[font=\tiny, align=center] at (0.25,5.5)
      {Narrow\\$r_3$};
  \end{scope}

  % Formula
  \node[codexblue, font=\small, align=center]
    at (4.5,-0.6)
    {$h = \dfrac{2\gamma\cos\theta}{\rho g r}$
    \quad $h \propto \dfrac{1}{r}$};

\end{tikzpicture}
\end{center}
\end{diagbox}

\begin{examplebox}[12.3]
\stepcounter{workedexample}
Water ($\gamma = 0.073\ \text{N\,m}^{-1}$,
$\rho = 1000\ \text{kg\,m}^{-3}$) rises in a glass
capillary tube of internal radius $0.30\ \text{mm}$.
The contact angle is $0°$.
($g = 10\ \text{m\,s}^{-2}$)\\
(a) Calculate the height of the water column.\\
(b) What radius tube would cause the water to rise
$5.0\ \text{cm}$?\\
(c) Explain why water rises in the wick of a
kerosene lantern.
\end{examplebox}

\begin{solutionbox}
\textbf{(a)} Height of rise:
\[
  h = \frac{2\gamma\cos\theta}{\rho g r}
  = \frac{2 \times 0.073 \times 1}{1000 \times 10
    \times 0.30\times10^{-3}}
  = \frac{0.146}{3.0}
  = 0.0487\ \text{m} \approx 4.9\ \text{cm}
\]

\textbf{(b)} Required radius:
\[
  r = \frac{2\gamma}{\rho g h}
  = \frac{2 \times 0.073}{1000 \times 10 \times 0.050}
  = \frac{0.146}{500}
  = 2.92\times10^{-4}\ \text{m}
  \approx 0.29\ \text{mm}
\]

\textbf{(c)} A wick consists of tightly-packed fibres
with very narrow spaces (capillaries) between them.
These narrow capillaries pull kerosene upward by
capillary action. The narrower the gaps between
fibres, the higher the kerosene can be drawn. This
capillary-fed fuel then burns at the top of the wick.
The same principle operates in soil --- water is
drawn from deeper, wetter layers toward the surface
through narrow soil pores.
\end{solutionbox}

% ============================================================
\section{Viscosity}
% ============================================================

\begin{definitionbox}[Viscosity]
\textbf{Viscosity} is the property of a fluid that
resists the relative motion between adjacent layers
of the fluid. It is the fluid's internal friction.

A highly viscous fluid (like honey or palm oil) flows
slowly and resists deformation strongly. A fluid of
low viscosity (like water or air) flows easily.

The coefficient of viscosity $\eta$ (eta) is defined
through Newton's law of viscosity:
\[
  F = \eta A \frac{dv}{dy}
\]
where $F$ is the tangential force between fluid layers,
$A$ is the area of contact, and $dv/dy$ is the
velocity gradient (rate of change of velocity with
distance perpendicular to flow).

SI unit of $\eta$: pascal-second (Pa\,s) or
N\,s\,m$^{-2}$.
\end{definitionbox}

\begin{diagbox}[Laminar Flow: Velocity Profile in a Viscous Fluid]
\begin{center}
\begin{tikzpicture}[scale=1.0, font=\small]

  % Pipe walls
  \fill[gray!40] (-0.2,-0.3) rectangle (10,0);
  \fill[gray!40] (-0.2,3.0) rectangle (10,3.4);
  \draw[thick] (-0.2,0) -- (10,0);
  \draw[thick] (-0.2,3.0) -- (10,3.0);

  % Fluid layers (laminar)
  \fill[codexblue!10] (-0.2,0) rectangle (10,3.0);

  % Velocity arrows (parabolic profile)
  % Wall: v = 0; centre: v = max
  \foreach \y/\vlen in {0.15/0.2, 0.5/0.8, 1.0/1.6,
    1.5/2.0, 2.0/1.6, 2.5/0.8, 2.85/0.2} {
    \draw[->, thick, codexblue]
      (2,\y) -- ({2+\vlen},\y);
  }

  % Layer labels
  \node[black, font=\tiny] at (0.5,0.05) [below] {Wall ($v=0$)};
  \node[black, font=\tiny] at (0.5,2.95) [above] {Wall ($v=0$)};
  \node[codexblue, font=\tiny] at (4.8,1.5) {Centre ($v_{max}$)};

  % Velocity profile curve
  \draw[codexred, very thick, dashed]
    (2,0) .. controls (2.2,1.5) .. (4,1.5)
    .. controls (2.2,1.5) .. (2,3.0);

  % Flow direction
  \draw[->, very thick, codexgold]
    (7.1,1.5) -- (8.5,1.5)
    node[right]{Flow};

  % Shear stress annotation
  \draw[<->, codexred, dashed]
    (5.8,0.5) -- (5.8,2.5);
  \node[right, codexred, font=\tiny, align=left] at (5.9,1.5)
    {velocity\\gradient\\$dv/dy$};

  % Labels
  \node[font=\small, align=center] at (4.5,-0.8)
    {Laminar (smooth) flow: layers slide past each other\\
    Velocity is zero at walls, maximum at centre};

\end{tikzpicture}
\end{center}
\end{diagbox}

\begin{table}[H]
\centering
\caption{Viscosity of Common Fluids at 20°C}
\label{tab:viscosity}
\begin{tabular}{ll}
\toprule
\textbf{Fluid} & \textbf{Viscosity (Pa\,s)}\\
\midrule
Air & $1.8\times10^{-5}$\\
Water & $1.0\times10^{-3}$\\
Palm oil & $\approx 0.04$--$0.08$\\
Glycerol & $1.5$\\
Honey & $2$--$10$\\
Tar (at 20°C) & $>10^5$\\
\bottomrule
\end{tabular}
\end{table}

\begin{keybox}[Temperature and Viscosity]
\begin{itemize}[noitemsep]
  \item \textbf{Liquids:} Viscosity \textit{decreases}
  with increasing temperature (molecules have more
  energy to overcome intermolecular attractions and
  flow past each other more easily).
  \item \textbf{Gases:} Viscosity \textit{increases}
  with temperature (more molecular collisions between
  layers as temperature rises).
\end{itemize}
This difference is physically important: engine oil
must maintain adequate viscosity at high engine
temperatures; thin oil at high temperature fails to
lubricate correctly.
\end{keybox}

% ============================================================
\section{Stokes' Law and Terminal Velocity}
% ============================================================

When a sphere moves through a viscous fluid, the fluid
exerts a drag force on it. For slow speeds (laminar
flow), this drag force was calculated by George Stokes
in 1851.

\begin{lawbox}[Stokes' Law]
The viscous drag force $F_d$ on a sphere of radius
$r$ moving at velocity $v$ through a fluid of
viscosity $\eta$:
\[
  F_d = 6\pi\eta r v
\]
This assumes \textbf{laminar flow} (smooth, layered)
at low speeds. At high speeds, flow becomes turbulent
and Stokes' law no longer applies.
\end{lawbox}

\begin{processbox}[Terminal Velocity in a Viscous Fluid]
A sphere of radius $r$, density $\rho_s$, falling
through a fluid of density $\rho_f$ and viscosity $\eta$
experiences three forces:

\begin{itemize}[noitemsep]
  \item Weight (downward):
  $W = \frac{4}{3}\pi r^3 \rho_s g$
  \item Upthrust (upward):
  $U = \frac{4}{3}\pi r^3 \rho_f g$
  \item Stokes' drag (upward):
  $F_d = 6\pi\eta r v$
\end{itemize}

At terminal velocity $v_T$, net force $= 0$:
\[
  W = U + F_d
\]
\[
  \frac{4}{3}\pi r^3 \rho_s g
  = \frac{4}{3}\pi r^3 \rho_f g + 6\pi\eta r v_T
\]
\[
  6\pi\eta r v_T
  = \frac{4}{3}\pi r^3 g(\rho_s - \rho_f)
\]
\[
  \boxed{v_T = \frac{2r^2(\rho_s - \rho_f)g}{9\eta}}
\]
\end{processbox}

\begin{diagbox}[Forces on a Sphere Falling Through Viscous Fluid]
\begin{center}
\begin{tikzpicture}[scale=1.0, font=\small]

  % Fluid background
  \fill[codexblue!10] (-1.0,0) rectangle (6,8);
  \draw[thick] (-1.0,0) rectangle (6,8);
  \node[codexblue, font=\itshape] at (2.5,7.5)
    {Viscous fluid ($\eta$, $\rho_f$)};

  % Sphere
  \fill[gray!60] (2.5,4.5) circle (0.7cm);
  \draw[thick] (2.5,4.5) circle (0.7cm);
  \node[white, font=\bfseries\small] at (2.5,4.5) {$m$};

  % Weight arrow (down)
  \draw[->, very thick, codexred, line width=2.5pt]
    (2.5,3.8) -- (2.5,2.2)
    node[right]{$W = \dfrac{4}{3}\pi r^3 \rho_s g$};

  % Upthrust arrow (up)
  \draw[->, very thick, codexgreen, line width=2pt]
    (2.5,5.2) -- (2.5,6.0)
    node[right]{$U = \dfrac{4}{3}\pi r^3 \rho_f g$};

  % Stokes drag arrow (up)
  \draw[->, very thick, codexblue, line width=2pt]
    (1.8,5.2) -- (1.8,6.0)
    node[left]{$F_d = 6\pi\eta r v$};

  % Velocity arrow (down, sphere moving)
  \draw[->, thick, codexgold, dashed]
    (3.2,4.5) -- (3.2,3.3)
    node[right, font=\small]{$v$ (velocity)};

  % At terminal velocity annotation
  \draw[thick, black, dashed]
    (0.2,1.5) -- (4.8,1.5);
  \node[font=\small, align=center] at (2.5,0.8)
    {At terminal velocity:\\
    $W = U + F_d$\quad (schematic arrows, not to scale; net force = 0)};

  % Terminal velocity formula
  \node[codexblue, font=\small, align=center]
    at (2.5,-0.8)
    {$v_T = \dfrac{2r^2(\rho_s-\rho_f)g}{9\eta}$};

\end{tikzpicture}
\end{center}
\end{diagbox}

\begin{examplebox}[12.4]
\stepcounter{workedexample}
A steel ball bearing of radius $2.0\ \text{mm}$ and
density $7800\ \text{kg\,m}^{-3}$ falls through a
cylinder of glycerol ($\eta = 1.5\ \text{Pa\,s}$,
$\rho_f = 1260\ \text{kg\,m}^{-3}$).\\
(a) Calculate the terminal velocity.\\
(b) Calculate the weight, upthrust and Stokes' drag
at terminal velocity.\\
(c) Verify that the three forces balance.
($g = 10\ \text{m\,s}^{-2}$)
\end{examplebox}

\begin{solutionbox}
\textbf{(a)} Terminal velocity:
\[
  v_T = \frac{2r^2(\rho_s-\rho_f)g}{9\eta}
  = \frac{2 \times (2.0\times10^{-3})^2
    \times (7800-1260) \times 10}
    {9 \times 1.5}
\]
\[
  = \frac{2 \times 4\times10^{-6} \times 6540 \times 10}
    {13.5}
  = \frac{5.232\times10^{-1}}{13.5}
  = 0.03876\ \text{m\,s}^{-1}
  \approx 3.9\ \text{cm\,s}^{-1}
\]

\textbf{(b)} Volume of sphere:
$V = \frac{4}{3}\pi r^3
= \frac{4}{3}\pi(2\times10^{-3})^3
= 3.351\times10^{-8}\ \text{m}^3$

Weight:
$W = \rho_s V g = 7800 \times 3.351\times10^{-8} \times 10
= 2.614\times10^{-3}\ \text{N}$

Upthrust:
$U = \rho_f V g = 1260 \times 3.351\times10^{-8} \times 10
= 4.222\times10^{-4}\ \text{N}$

Stokes' drag:
$F_d = 6\pi\eta r v_T
= 6\pi \times 1.5 \times 2\times10^{-3} \times 0.03876
= 2.192\times10^{-3}\ \text{N}$

\textbf{(c)} Verification:
$U + F_d = 4.222\times10^{-4} + 2.192\times10^{-3}
= 2.614\times10^{-3}\ \text{N} = W$ ✓
\end{solutionbox}

% ============================================================
\section{Laminar vs Turbulent Flow}
% ============================================================

\begin{definitionbox}[Laminar and Turbulent Flow]
\textbf{Laminar flow:} Fluid moves in smooth,
parallel layers (streamlines) with no mixing between
layers. Velocity is steady at each point. Occurs at
low speeds and in viscous fluids.

\vspace{0.4cm}

\textbf{Turbulent flow:} Fluid moves in irregular,
chaotic eddies. Layers mix. Velocity fluctuates at
each point. Occurs at high speeds and in less viscous
fluids.

The transition is predicted by the
\textbf{Reynolds number}:
\[
  Re = \frac{\rho v d}{\eta}
\]
where $d$ is a characteristic length (e.g.\ pipe
diameter). Laminar flow: $Re < 2000$. Turbulent:
$Re > 4000$. Transition: $2000 < Re < 4000$.
\end{definitionbox}

\begin{diagbox}[Laminar vs Turbulent Flow]
\begin{center}
\begin{tikzpicture}[scale=0.9, font=\small]

  % === LAMINAR FLOW ===
  \begin{scope}[shift={(0,0)}]
    \node[font=\bfseries] at (3,5.5) {Laminar Flow};
    % Pipe
    \fill[gray!20] (0,-0.3) rectangle (6,0);
    \fill[gray!20] (0,3.0) rectangle (6,3.3);
    \draw[thick] (0,0) -- (6,0);
    \draw[thick] (0,3.0) -- (6,3.0);
    \fill[codexblue!10] (0,0) rectangle (6,3.0);

    % Smooth streamlines
    \foreach \y in {0.3,0.8,1.5,2.2,2.7} {
      \draw[codexblue, thick]
        (0.3,\y) -- (5.7,\y);
    }
    % Flow arrows
    \foreach \y in {0.3,1.5,2.7} {
      \draw[->, very thick, codexblue]
        (2.5,\y) -- (3.5,\y);
    }
    \node[font=\small, codexblue, align=center]
      at (3,4.5) {Smooth, parallel\\streamlines};
  \end{scope}

  % === TURBULENT FLOW ===
  \begin{scope}[shift={(7,0)}]
    \node[font=\bfseries] at (3,5.5) {Turbulent Flow};
    % Pipe
    \fill[gray!20] (0,-0.3) rectangle (6,0);
    \fill[gray!20] (0,3.0) rectangle (6,3.3);
    \draw[thick] (0,0) -- (6,0);
    \draw[thick] (0,3.0) -- (6,3.0);
    \fill[codexblue!10] (0,0) rectangle (6,3.0);

    % Chaotic streamlines
    \draw[codexblue, thick]
      (0.3,1.5) .. controls (1,0.5) and (1.5,2.5)
      .. (2.5,1.5) .. controls (3.5,0.5) and (4,2.5)
      .. (5,1.5) -- (5.7,1.5);
    \draw[codexblue, thick]
      (0.3,0.8) .. controls (1,1.5) and (1.5,0.3)
      .. (2.5,0.8) .. controls (3,1.5) and (4,0.3)
      .. (5,1.0) -- (5.7,1.2);
    \draw[codexblue, thick]
      (0.3,2.2) .. controls (1,1.2) and (1.5,2.8)
      .. (2.5,2.2) .. controls (3.5,1.5) and (4,2.7)
      .. (5,2.0) -- (5.7,2.1);

    % Eddy arrows
    \draw[->, codexred, thick]
      (2.5,1.5) circle (0.4);
    \draw[->, codexred, thick]
      (4.0,1.2) circle (0.35);

    \node[font=\small, codexred, align=center]
      at (3,4.5) {Chaotic eddies,\\mixing, energy loss};
  \end{scope}

\end{tikzpicture}
\end{center}
\end{diagbox}

% ============================================================
\section{Chapter Summary}
% ============================================================

\begin{summarybox}
\textbf{Surface tension origin:} Net inward cohesive
force on surface molecules. Surface contracts to
minimise area.

\textbf{Surface tension:}
$\gamma = F/L = W/A$ (N\,m$^{-1}$)

\textbf{Excess pressure:}
Liquid drop: $\Delta P = 2\gamma/r$
(one surface)\\
Soap bubble: $\Delta P = 4\gamma/r$
(two surfaces)

\textbf{Capillarity:} Water rises in narrow glass
tubes (adhesion $>$ cohesion, concave meniscus).
Mercury falls (cohesion $>$ adhesion, convex meniscus).

\textbf{Capillary rise:}
$h = 2\gamma\cos\theta/(\rho g r)$;
narrower tube $\implies$ greater rise.

\textbf{Viscosity:} Internal friction of a fluid.
Liquids: $\eta$ decreases with temperature.
Gases: $\eta$ increases with temperature.

\textbf{Stokes' law:} $F_d = 6\pi\eta rv$
(sphere in laminar flow)

\textbf{Terminal velocity:}
$v_T = 2r^2(\rho_s-\rho_f)g/(9\eta)$

\textbf{Laminar flow:} smooth, parallel layers,
$Re < 2000$.
\textbf{Turbulent flow:} chaotic, mixing,
$Re > 4000$.
\end{summarybox}

% ============================================================
\newpage
\begin{exercisebox}[12]

\subsection*{Section A: Multiple Choice}

\textbf{A1.} Surface tension arises because molecules
at the liquid surface:

\mcoption{A}{Have more kinetic energy than interior
molecules}
\mcoption{B}{Experience a net inward cohesive force}
\mcoption{C}{Are repelled by air molecules above
them}
\mcoption{D}{Have lower density than interior
molecules}
\ansref{Ch12-A1}

\vspace{0.4cm}

\textbf{A2.} The excess pressure inside a soap bubble
of radius $r$ and surface tension $\gamma$ is:

\mcoption{A}{$\gamma/r$}
\mcoption{B}{$2\gamma/r$}
\mcoption{C}{$4\gamma/r$}
\mcoption{D}{$8\gamma/r$}
\ansref{Ch12-A2}

\vspace{0.4cm}

\textbf{A3.} Water rises in a capillary tube of radius
$r$ to height $h$. In a tube of radius $r/3$, the
height of rise would be:

\mcoption{A}{$h/3$}
\mcoption{B}{$h/9$}
\mcoption{C}{$3h$}
\mcoption{D}{$9h$}
\ansref{Ch12-A3}

\vspace{0.4cm}

\textbf{A4.} Mercury forms a convex meniscus in a
glass tube because:

\mcoption{A}{Mercury is denser than glass}
\mcoption{B}{Mercury's cohesion exceeds its adhesion
to glass}
\mcoption{C}{Mercury is a liquid metal}
\mcoption{D}{Glass repels mercury molecules}
\ansref{Ch12-A4}

\vspace{0.4cm}

\textbf{A5.} Adding soap to water:

\mcoption{A}{Increases surface tension}
\mcoption{B}{Has no effect on surface tension}
\mcoption{C}{Reduces surface tension}
\mcoption{D}{Reduces viscosity only}
\ansref{Ch12-A5}

\vspace{0.4cm}

\textbf{A6.} Stokes' law $F = 6\pi\eta rv$ applies
when:

\mcoption{A}{The flow around the sphere is turbulent}
\mcoption{B}{The sphere is hollow}
\mcoption{C}{The flow around the sphere is laminar}
\mcoption{D}{The fluid is a gas only}
\ansref{Ch12-A6}

\vspace{0.4cm}

\textbf{A7.} A sphere reaches terminal velocity in
a viscous fluid when:

\mcoption{A}{Gravity equals upthrust}
\mcoption{B}{Gravity equals Stokes' drag}
\mcoption{C}{Gravity equals upthrust plus
Stokes' drag}
\mcoption{D}{Stokes' drag equals upthrust}
\ansref{Ch12-A7}

\vspace{0.4cm}

\textbf{A8.} The terminal velocity of a sphere is
proportional to:

\mcoption{A}{$r$}
\mcoption{B}{$r^2$}
\mcoption{C}{$r^3$}
\mcoption{D}{$1/r$}
\ansref{Ch12-A8}

\vspace{0.4cm}

\textbf{A9.} The viscosity of a liquid generally:

\mcoption{A}{Increases with temperature}
\mcoption{B}{Decreases with temperature}
\mcoption{C}{Is independent of temperature}
\mcoption{D}{Increases then decreases with temperature}
\ansref{Ch12-A9}

\vspace{0.4cm}

\textbf{A10.} Two soap bubbles of radii $r$ and $3r$
are connected by a tube. Air will flow:

\mcoption{A}{From the larger bubble to the smaller}
\mcoption{B}{From the smaller bubble to the larger}
\mcoption{C}{In both directions equally}
\mcoption{D}{No air flows; the system is already
in equilibrium}
\ansref{Ch12-A10}

\vspace{0.6cm}

\subsection*{Section B: Theory and Calculations}

\textbf{B1.} A thin wire loop of circumference
$12\ \text{cm}$ is placed flat on a water surface.
The force required to pull the loop vertically out
of the water (breaking the surface film) is
$1.75\times10^{-2}\ \text{N}$.\\[4pt]
(a) Explain why the factor of 2 appears in the
surface tension calculation for a wire loop pulled
out of water (think about how many surfaces are
created).\\
(b) Calculate the surface tension of water.\\
(c) What additional weight could be supported on the
surface of water by a straight wire of length
$5.0\ \text{cm}$? (Treat the wire as resting on the
surface, supported by the film on both sides.)
\hfill\ansref{Ch12-B1}

\vspace{0.4cm}

\textbf{B2.} A capillary tube of internal radius
$0.20\ \text{mm}$ stands vertically in a liquid of
density $800\ \text{kg\,m}^{-3}$ and surface tension
$0.050\ \text{N\,m}^{-1}$. The contact angle between
the liquid and glass is $30°$.
($g = 10\ \text{m\,s}^{-2}$, $\cos 30° = 0.866$)\\[4pt]
(a) Calculate the height to which the liquid rises.\\
(b) If the tube is only $5.0\ \text{cm}$ tall, the
liquid cannot rise to its natural height. Explain
what happens instead, using the capillary formula
to find the new contact angle.\\
(c) If the tube radius is doubled to $0.40\ \text{mm}$,
what is the new height of rise?\\
(d) Explain the physical significance of capillarity
for water transport in trees.
\hfill\ansref{Ch12-B2}

\vspace{0.4cm}

\textbf{B3.} A glass sphere of radius $3.0\ \text{mm}$
and density $2500\ \text{kg\,m}^{-3}$ falls through
castor oil of viscosity $0.985\ \text{Pa\,s}$ and
density $960\ \text{kg\,m}^{-3}$.
($g = 10\ \text{m\,s}^{-2}$)\\[4pt]
(a) Calculate the terminal velocity of the sphere.\\
(b) Sketch a velocity-time graph from the moment the
sphere is released to when it reaches terminal velocity.
Label the key features.\\
(c) At the instant of release ($v = 0$), calculate
the initial acceleration.\\
(d) If the sphere were twice the radius but the same
density, how would the terminal velocity change?
Show the calculation.
\hfill\ansref{Ch12-B3}

\vspace{0.4cm}

\textbf{B4.} A soap bubble is blown to a radius of
$4.0\ \text{cm}$. The surface tension of the soap
solution is $0.040\ \text{N\,m}^{-1}$.
Atmospheric pressure is $1.013\times10^5\ \text{Pa}$.\\[4pt]
(a) Calculate the excess pressure inside the bubble.\\
(b) Calculate the total pressure inside.\\
(c) The bubble is blown further until its radius
doubles to $8.0\ \text{cm}$. Calculate the work done
in increasing the surface area. (Use $W = \gamma\Delta A$,
remembering two surfaces.)\\
(d) A second bubble of radius $2.0\ \text{cm}$ is
connected to the first. Predict and calculate the
final radius of the combined bubble, using conservation
of air volume (assume the bubbles merge into one).
\hfill\ansref{Ch12-B4}

\vspace{0.4cm}

\textbf{B5.} An engineer uses a falling-sphere
viscometer to measure the viscosity of an unknown oil.
A steel ball bearing of radius $1.5\ \text{mm}$
and density $7800\ \text{kg\,m}^{-3}$ is dropped into
the oil (density $870\ \text{kg\,m}^{-3}$).
The ball is observed to fall $20\ \text{cm}$ in
$8.5\ \text{s}$ at constant speed.
($g = 10\ \text{m\,s}^{-2}$)\\[4pt]
(a) Calculate the terminal velocity of the ball.\\
(b) Use Stokes' law to calculate the viscosity of
the oil.\\
(c) Calculate the Reynolds number using
$Re = \rho_f v d/\eta$. Is Stokes' law valid in
this case?\\
(d) At what temperature do you expect the oil to
be? (Refer to Table~\ref{tab:viscosity} and
physical reasoning.)\\
(e) If the experiment is repeated with a ball of
twice the radius (keeping everything else the same),
what is the new terminal velocity? What does this
imply about using larger spheres in viscometry?
\hfill\ansref{Ch12-B5}

\end{exercisebox}
